%===================================================================
% Chapter 7: Conclusion
%===================================================================
\chapter{Conclusion}
\label{chap:conclusion}

For four decades, automatic query rewriting has meant applying rules
that someone wrote down in advance. This dissertation asked whether
that boundary is necessary, and showed that it is not: semantically
equivalent, substantially faster SQL can be discovered rather than
pattern-matched, for individual queries and for the pipelines in
which modern analytics runs.

The path ran along two movements. The source of rewrites moved from
systematic search to learned knowledge: SlabCity established that
whole-query rewriting needs no rules at all, and GenRewrite showed
that language models, properly disciplined, optimize the complex
queries where search falters and grow more capable with every query
they optimize. The object of rewriting moved from the single query
to the pipeline: DAGSmith treated dependency structure as
optimization signal and reached improvements that no
one-query-at-a-time technique can see. What never changed is the
discipline underneath. Every candidate, from whatever source, earned
its place by passing equivalence checking matched to its setting.
Removing the rules did not remove the guarantees; it moved them from
the transformation to the result.

The trends that motivated this work have only strengthened. SQL is
written increasingly by AI, accumulates increasingly into pipelines,
and runs on infrastructure that bills inefficiency by the second,
so the space of queries that no hand-written rule set will
anticipate keeps growing. The central lesson of this dissertation is
that this space is not a wilderness: it can be searched, it can be
learned, and it can be verified. The rules we can write down were
never the boundary of what optimization can reach; they were only
where it began.